\section{Case Study 3: Dokumentation der Durchführung}
\label{appendix:case_study_3_documentation}

\subsection{Problem Statement - Risiken} \label{appendix:cs3_phase1_risiken}
\begin{table}[H]
\centering
\renewcommand{\arraystretch}{1.3}
\begin{tabular}{|p{3cm}|p{4.9cm}|p{3cm}|p{4.9cm}|p{4.9cm}|}
\hline
\textbf{Risiko} & \textbf{Eintrittswahrscheinlichkeit} & \textbf{Auswirkung} & \textbf{Entdeckbarkeit} & \textbf{KRI} \\
\hline

\end{tabular}
\caption{Risikobewertung der Phase 1 nach erweitertem FMEA Schema}
\label{tab:cs3_phase1_risiken_tabelle}
\end{table}

\subsection{Problem Statement - KRI nachverfolgung} \label{appendix:cs3_phase1_risiko_verfolgung}
\begin{table}[H]
\centering
\renewcommand{\arraystretch}{1.3}
\begin{tabular}{|p{3cm}|p{4.9cm}|p{3cm}|p{4.9cm}|}
\hline
\textbf{Risiko} & \textbf{KRI vorher} & \textbf{Maßnahme} & \textbf{KRI nachher} \\
\hline

\end{tabular}
\caption{Risikoregister: Kriterien zur Bildung der RPZ innerhalb des 1. Risk Gates}

\label{tab:cs3_phase1_risiko_verfolgung_tabelle}
\end{table}

% --------------------------- Data Provisioning ------------------------------%

\subsection{Data Provisioning - Risiken} \label{appendix:cs3_phase2_risiken}
\begin{table}[H]
\centering
\renewcommand{\arraystretch}{1.3}
\begin{tabular}{|p{3cm}|p{4.9cm}|p{3cm}|p{4.9cm}|p{4.9cm}|}
\hline
\textbf{Risiko} & \textbf{Eintrittswahrscheinlichkeit} & \textbf{Auswirkung} & \textbf{Entdeckbarkeit} & \textbf{KRI} \\
\hline

Datenqualität zu gering
& 2 
& 5
& 2
& 3.2 \\
\hline

Infrastruktur nicht ausreichend 
& 2 
& 5 
& 2 
& 3.2 \\
\hline

\end{tabular}
\caption{Risikobewertung der Phase 2 nach erweitertem FMEA Schema}
\label{tab:cs3_phase2_risiken_tabelle}
\end{table}

\subsection{Data Provisioning - KRI nachverfolgung} \label{appendix:cs3_phase2_risiko_verfolgung}
\begin{table}[H]
\centering
\renewcommand{\arraystretch}{1.3}
\begin{tabular}{|p{3cm}|p{4.9cm}|p{3cm}|p{4.9cm}|}
\hline
\textbf{Risiko} & \textbf{KRI vorher} & \textbf{Maßnahme} & \textbf{KRI nachher} \\
\hline

Datenqualität zu gering
& 0
& Beobachtung der Performance
& 3.2 \\
\hline

Infrastruktur nicht ausreichend 
& 0
& Beobachtung der Infrastruktur
& 3.2 \\
\hline

\end{tabular}
\caption{Risikoregister: Kriterien zur Bildung der RPZ innerhalb des 2. Risk Gates}

\label{tab:cs3_phase2_risiko_verfolgung_tabelle}
\end{table}

% --------------------------- Analysis ------------------------------%

\subsection{Analysis - Risiken im ersten Anlauf} \label{appendix:cs3_kri_mapping_1}
\begin{table}[H]
\centering
\renewcommand{\arraystretch}{1.3}
\begin{tabular}{p{5cm} c p{4cm} c c c c}
\hline
\textbf{Risiko (inkl. Interpretation)} & \textbf{Metrikwert} & \textbf{Mapping} & \textbf{Eintrittswahrscheinlichkeit} & \textbf{Auswirkung} & \textbf{Entdeckbarkeit} & \textbf{KRI} \\
\hline

Fehlklassifikation: Nicht betäubter Fisch wird als betäubt erkannt (kritischer False Negative Fall) 
& 25\% (FNR) 
& $<2\%$ = gering; $>2\%$ = mittel; $>$10\% = hoch 
& 4 
& 5
& 2 
& 4 \\

Unzureichende Erkennungsrate tatsächlich betäubter Fische (Recall-Schwäche)
& 72\% (Recall) 
& 0--90\% = kritisch; 90 -- 95\% = ausreichend; > 95\% = sehr gut
& 4
& 5 
& 2 
& 4 \\

Geringe Trennschärfe des Modells (unsichere Klassifikation)
& 0{,}80 (ROC-AUC) 
& 0 -- 0{,}74 = schlecht , 0{,}75--0{,}85 = mittel, 0{,}86--1{,}00 = gut 
& 4 
& 5 
& 2 
& 4 \\

\hline
\end{tabular}
\caption{Überführung von Modellperformanzmetriken in risikobezogene Kennziffern (KRI)}
\label{tab:cs3_phase3_kri_mapping_tabelle_1}
\end{table}

\subsection{Analysis - KRI Nachverfolgung im ersten Anlauf} \label{appendix:cs3_phase3_risiko_verfolgung_1}
\begin{table}[H]
\centering
\renewcommand{\arraystretch}{1.3}
\begin{tabular}{|p{3cm}|p{4.9cm}|p{3cm}|p{4.9cm}|}
\hline
\textbf{Risiko} & \textbf{KRI vorher} & \textbf{Maßnahme} & \textbf{KRI nachher} \\
\hline

Datenqualität zu gering
& 3.2
& Beobachtung der Performance
& 3.2 \\
\hline

Infrastruktur nicht ausreichend 
& 3.2
& Beobachtung der Infrastruktur
& 3.2 \\
\hline

Fehlklassifikation: Nicht betäubter Fisch wird als betäubt erkannt (kritischer False Negative Fall) 
& 0
& Upgrade der Hardware
& 4  \\
\hline

Unzureichende Erkennungsrate tatsächlich betäubter Fische (Recall-Schwäche)
& 0
& Upgrade der Hardware
& 4  \\
\hline

Geringe Trennschärfe des Modells (unsichere Klassifikation)
& 0
& Upgrade der Hardware
& 4  \\
\hline

\end{tabular}
\caption{Risikoregister: Kriterien zur Bildung der RPZ innerhalb des 3. Risk Gates - Erster Anlauf}

\label{tab:cs3_phase3_risiko_verfolgung_tabelle_1}
\end{table}


\subsection{Analysis - Risiken im zweiten Anlauf} \label{appendix:cs3_kri_mapping_2}
\begin{table}[H]
\centering
\renewcommand{\arraystretch}{1.3}
\begin{tabular}{p{5cm} c p{4cm} c c c c}
\hline
\textbf{Risiko (inkl. Interpretation)} & \textbf{Metrikwert} & \textbf{Mapping} & \textbf{Eintrittswahrscheinlichkeit} & \textbf{Auswirkung} & \textbf{Entdeckbarkeit} & \textbf{KRI} \\
\hline

Akzeptable Fehlklassifikationswahrscheinlichkeit: Nicht betäubter Fisch wird als betäubt erkannt (kritischer False Negative Fall) 
& 9\% (FNR) 
& $<2\%$ = gering; $>2\%$ = mittel; $>$10\% = hoch 
& 3 
& 5 
& 2 
& 3.6 \\

Ausreichende Erkennungsrate tatsächlich betäubter Fische (Recall-Schwäche)
& 90\% (Recall) 
& 0--90\% = kritisch; 90 -- 95\% = ausreichend; > 95\% = sehr gut
& 3 
& 5 
& 2 
& 3.6 \\

Ausreichende Trennschärfe des Modells (unsichere Klassifikation)
& 0{,}90 (ROC-AUC) 
& 0 -- 0{,}74 = schlecht , 0{,}75--0{,}85 = mittel, 0{,}86--1{,}00 = gut
& 3 
& 5 
& 2 
& 3.6 \\

Datenqualität zu gering
& - 
& -
& 1 
& 5 
& 1
& 2.6 \\
\hline

\hline
\end{tabular}
\caption{Überführung von Modellperformanzmetriken in risikobezogene Kennziffern (KRI)}
\label{tab:cs3_phase3_kri_mapping_tabelle_2}
\end{table}

\subsection{Analysis - KRI Nachverfolgung im zweiten Anlauf} \label{appendix:cs3_phase3_risiko_verfolgung_2}
\begin{table}[H]
\centering
\renewcommand{\arraystretch}{1.3}
\begin{tabular}{|p{3cm}|p{4.9cm}|p{3cm}|p{4.9cm}|}
\hline
\textbf{Risiko} & \textbf{KRI vorher} & \textbf{Maßnahme} & \textbf{KRI nachher} \\
\hline

Datenqualität zu gering
& 3.2
& Beobachtung der Performance
& 2.6 \\
\hline

Infrastruktur nicht ausreichend 
& 3.2
& Beobachtung der Infrastruktur
& 3.2 \\
\hline

Fehlklassifikation: Nicht betäubter Fisch wird als betäubt erkannt (kritischer False Negative Fall) 
& 4
& Weiterhin zur Beobachtung
& 3.6  \\
\hline

Unzureichende Erkennungsrate tatsächlich betäubter Fische (Recall-Schwäche)
& 4
& Weiterhin zur Beobachtung
& 3.6  \\
\hline

Geringe Trennschärfe des Modells (unsichere Klassifikation)
& 4
& Weiterhin zur Beobachtung
& 3.6  \\
\hline

\end{tabular}
\caption{Risikoregister: Kriterien zur Bildung der RPZ innerhalb des 3. Risk Gates - Zweiter Anlauf}

\label{tab:cs3_phase3_risiko_verfolgung_tabelle_2}
\end{table}

% --------------------------- DEPLOYMENT ------------------------------%


\subsection{Deployment - Risiken} \label{appendix:cs3_phase4_risiken}
\begin{table}[H]
\centering
\renewcommand{\arraystretch}{1.3}
\begin{tabular}{p{5cm} c p{4cm} c c c c}
\hline
\textbf{Risiko (inkl. Interpretation)} & \textbf{Metrikwert} & \textbf{Mapping} & \textbf{Eintrittswahrscheinlichkeit} & \textbf{Auswirkung} & \textbf{Entdeckbarkeit} & \textbf{KRI} \\
\hline

Instabilität bei veränderten Umweltbedingungen (z.\,B. Lichtverhältnisse, Data Drift)
& 5\% Performance Drop 
& 0 -- 5\% = akzeptabel; 5 -- 10\% = mittel ; 10--20\% = kritisch 
& 2 
& 4 
& 2 
& 2.8 \\

\hline
\end{tabular}
\caption{Überführung von Modellperformanzmetriken in risikobezogene Kennziffern (KRI)}
\label{tab:cs3_phase4_kri_mapping}
\end{table}

\subsection{Deployment - KRI nachverfolgung} \label{appendix:cs3_phase4_risiko_verfolgung}
\begin{table}[H]
\centering
\renewcommand{\arraystretch}{1.3}
\begin{tabular}{|p{3cm}|p{4.9cm}|p{3cm}|p{4.9cm}|}
\hline
\textbf{Risiko} & \textbf{KRI vorher} & \textbf{Maßnahme} & \textbf{KRI nachher} \\
\hline

Datenqualität zu gering
& 3.2
& Beobachtung der Performance
& 2.6 \\
\hline

Infrastruktur nicht ausreichend 
& 3.2
& Beobachtung der Infrastruktur
& 3.2 \\
\hline

Fehlklassifikation: Nicht betäubter Fisch wird als betäubt erkannt (kritischer False Negative Fall) 
& 4
& Weiterhin zur Beobachtung
& 3.6  \\
\hline

Unzureichende Erkennungsrate tatsächlich betäubter Fische (Recall-Schwäche)
& 4
& Weiterhin zur Beobachtung
& 3.6  \\
\hline

Geringe Trennschärfe des Modells (unsichere Klassifikation)
& 4
& Weiterhin zur Beobachtung
& 3.6  \\
\hline

Instabilität bei veränderten Umweltbedingungen (z.\,B. Lichtverhältnisse, Data Drift)
& 0
& Beobachtung. Bei Verschlechterung eine Beleuchtungslösung anbauen.
& 2.8 \\
\hline

\end{tabular}
\caption{Risikoregister: Kriterien zur Bildung der RPZ innerhalb des 4. Risk Gates}

\label{tab:cs3_phase4_risiko_verfolgung_tabelle}
\end{table}
